\section{Machine-checked deterministic theorem and its correspondence}
\label{app:lean-formalization}

This appendix records the public theorem exported by the Lean~4
development~\cite{DeMouraUllrich2021} associated with this manuscript and
explains how its formal predicates correspond to the deterministic statements
in \cref{thm:det-obligatory,thm:optional-curve}.  The development uses the
mathlib library~\cite{Mathlib2020} and is pinned to Lean and mathlib
version~4.30.0.  Its root module imports
\texttt{SchedulingPaper.MainTheorems}, whose final theorem is reproduced
below.  Apart from the typographic rendering of Unicode symbols, this is a
verbatim excerpt.

\begin{lstlisting}[style=lean]
theorem exact_ratio_main
    : (∀ {u : ℝ}, 0 < u →
      FiniteExactRatioConclusion u (exactCurve u)) ∧
      ObligatoryExactRatioConclusion RStar := by
  constructor
  · intro u hu
    exact finite_exact_ratio hu
  · exact obligatory_exact_ratio
\end{lstlisting}

The two conclusion records occurring in the theorem are:

\begin{lstlisting}[style=lean]
structure FiniteExactRatioConclusion (u ratio : ℝ) : Prop where
  lower : LowerBound.FixedSizeLowerBound (.finite u) ratio
  upper :
    ∃ strategy : ∀ n, Online.Strategy n,
      UpperBound.CompletesByAnalysisFuel (.finite u) strategy ∧
      SizeAsymptoticUpper
        (UpperBound.strategyCost (.finite u) strategy)
        (UpperBound.fixedOfflineCost (.finite u)) ratio

structure ObligatoryExactRatioConclusion (ratio : ℝ) : Prop where
  lower : LowerBound.FixedSizeLowerBound .infinite ratio
  upper :
    ∃ strategy : ∀ n, Online.Strategy n,
      UpperBound.CompletesByAnalysisFuel .infinite strategy ∧
      SizeAsymptoticUpper
        (UpperBound.strategyCost .infinite strategy)
        (UpperBound.fixedOfflineCost .infinite) ratio
\end{lstlisting}

\subsection{Objects on the two sides}

The following table gives the direct dictionary between the formal objects
and the notation of the manuscript.

\begin{center}
\small
\begin{tabular}{@{}p{0.25\textwidth}p{0.29\textwidth}p{0.36\textwidth}@{}}
\toprule
Lean object & Manuscript object & Meaning \\
\midrule
\texttt{Cap.finite u}
  & finite common upper limit \(u\)
  & processing times lie in \([0,u]\); testing and raw execution are
    available \\
\texttt{Cap.infinite}
  & obligatory testing
  & processing times are nonnegative and raw execution is unavailable \\
\texttt{exactCurve u}
  & \(\V(u)\)
  & the six branches in \cref{eq:opt:full-curve}, including all five
    verified joins \\
\texttt{RStar}
  & \(\Rstar\)
  & \(1+2/(\zstar-1)\), where \(\zstar>1\) is the unique solution of
    \(\zstar-3=\log\zstar\) \\
\texttt{fixedOfflineCost}
  & \(\OPT\)
  & the clairvoyant shortest-first cost of the effective lengths
    \(q_u(p)\) \\
\texttt{strategyCost}
  & \(\ALG\)
  & the total completion cost of the operational online transcript \\
\bottomrule
\end{tabular}
\end{center}

A formal strategy on \(n\) jobs is a deterministic function from the public
transcript to its next action.  A test publishes the revealed processing
time; raw execution publishes no hidden value.  Thus the strategy has exactly
the information available to an online algorithm in the manuscript.  The
formal action type uses the standard no-idling normalization appropriate when
all jobs are released at time zero.  Processing times and comparisons are
exact real numbers.

\subsection{Expansion of the lower and upper fields}

For a cap and a candidate value \(R\),
\texttt{FixedSizeLowerBound} expands to the following statement.  For every
size-indexed family of deterministic strategies \(\mathcal B_n\), every
\(\eps>0\), and every \(N\), there are an \(n\ge N\), a genuine fixed
admissible \(n\)-job input \(I\), and a finite settled execution such that
either the execution leaves a job incomplete or
\begin{equation}
  (R-\eps)\OPT(I)\le \ALG_{\mathcal B_n}(I).
  \label{eq:lean-lower-expansion}
\end{equation}
An incomplete settled run represents infinite cost.  In particular, for
algorithms that complete all inputs, \eqref{eq:lean-lower-expansion} is the
usual lower-bound sequence on arbitrarily large fixed instances.  The lower
constructions are first described by adaptive revelation oracles, but the
formal replay theorem freezes every such interaction into the fixed input
quantified here; this is the machine-checked counterpart of
\cref{lem:replay}.

The predicate \texttt{SizeAsymptoticUpper} is the epsilon formulation
\begin{equation}
 \forall\eps>0\ \exists N\ \forall n\ge N\ \forall I,\qquad
 \ALG_{\mathcal A_n}(I)\le (R+\eps)\OPT(I).
 \label{eq:lean-upper-expansion}
\end{equation}
The accompanying completion field proves that every job is done after the
common analysis fuel \(2n+1\), rather than merely assigning a cost to a
partial trace.  Consequently,
\cref{eq:lean-lower-expansion,eq:lean-upper-expansion} are the two matching
operational statements that characterize the optimal value in
\cref{eq:cr-infty}.  The Lean theorem packages these two directions; it does
not introduce a separate real-valued infimum or \texttt{limsup} object.

Before projecting to \texttt{SizeAsymptoticUpper}, the checked upper proof
constructs a stronger \texttt{FixedUpperCertificate}: for the chosen cap it
contains a constant \(C\ge0\) such that, for every \(n\) and every fixed
input,
\begin{equation}
  \ALG_{\mathcal A_n}(I)\le R\OPT(I)+Cn.
  \label{eq:lean-linear-certificate}
\end{equation}
The quadratic lower bounds
\cref{eq:opt-n2,eq:optional-opt-n2} then imply
\eqref{eq:lean-upper-expansion}.  For \(u\ge\zstar\), the separately verified
plateau certificate uses the same strategy and the same \(C\), independent
of \(u\), including the obligatory endpoint.  This uniformity is stronger
than the quantifier order displayed by \texttt{exact\_ratio\_main}, which
only exposes one strategy after a cap has been fixed.

\subsection{Correspondence with the two deterministic theorems}

The first conjunct of \texttt{exact\_ratio\_main} quantifies over every
finite \(u>0\).  Substituting \(R=\texttt{exactCurve}\ u=\V(u)\), its lower
field gives the fixed-instance lower direction of
\cref{thm:optional-curve}, while its upper field supplies a completing
strategy and the matching upper direction.  The definition of
\texttt{exactCurve} is literally the six-way case split in
\cref{eq:opt:full-curve}; the Lean development also proves its joins,
continuity, plateau, and global maximum at \(\udiam\).

The second conjunct substitutes \(R=\texttt{RStar}=\Rstar\) in the
obligatory model.  Its lower field gives the lower half of
\cref{thm:det-obligatory}; its completing upper strategy together with
\eqref{eq:lean-linear-certificate} gives the upper half.  The common formal
infrastructure also
proves that the offline benchmark is the shortest-first schedule of
effective lengths and that the transcript cost is exactly its
suffix-weighted pair-accounting expression.

\subsection{Proof boundary}

The root Lean module imports the complete lower- and upper-bound assemblies.
The public theorem has no operational interface hypotheses, and the
development contains no \texttt{sorry}, \texttt{admit}, or project-defined
\texttt{axiom}.  It does use mathlib's classical real analysis and
noncomputable choices of exact real parameters.  Thus it certifies
deterministic mathematical strategy terms, not an extracted floating-point
implementation.

The theorem recorded here is specifically the size-asymptotic result.  It
does not assert the pointwise inequality \(\ALG\le R\OPT\) without a
remainder, nor admissibility of the endpoint coefficient with one fixed
additive constant.  The formal library separately derives additive
admissibility of every coefficient \(R+\eps\) from
\eqref{eq:lean-linear-certificate}.  The machine-checked mixed-regime upper
bound follows the cap-reserve potential proof; the bubbling argument in
\cref{app:cap-bubbling} remains an independent manuscript cross-check.
